% Part: many-valued-logic
% Chapter: infinite-valued-logics
% Section: lukasiewicz

\documentclass[../../../include/open-logic-section]{subfiles}

\begin{document}

\olfileid{mvl}{inf}{luk}

\olsection{\L ukasiewicz যুক্তিবিদ্যা}

\begin{editorial}
  এটি অসীমমানী \L ukasiewicz যুক্তিবিদ্যার একটি সংক্ষিপ্ত প্রাথমিক খসড়া।
\end{editorial}

\begin{defn}\ollabel{def:lukasiewicz} অসীমমানী \L ukasiewicz
যুক্তিবিদ্যা~$\LogLuk[\infty]$ নিচের ম্যাট্রিক্স দিয়ে সংজ্ঞায়িত:
\begin{enumerate}
  \item $\lnot$, $\land$, $\lor$, $\lif$-সহ প্রমিত বচনগত ভাষা
  $\Lang L_0$।
  \item সত্যমানের সেট $V_\infty$।
  \item $1$ একমাত্র মনোনীত মান; অর্থাৎ, $V^+ = \{1\}$।
  \item সত্যমান-অপেক্ষকগুলি নিচের অপেক্ষকগুলির মাধ্যমে নির্ধারিত:
  \begin{align*}
    \tf{\lnot}[\LogLuk](x) & = 1 - x\\
    \tf{\land}[\LogLuk](x,y) & = \min(x,y)\\
    \tf{\lor}[\LogLuk](x,y) & = \max(x,y)\\
    \tf{\lif}[\LogLuk](x,y) & = \min(1,1-(x-y)) = \begin{cases}
      1 & \text{if } x \le y\\
      1-(x-y) & \text{otherwise.}
    \end{cases}
    \end{align*}
\end{enumerate}
$m$-মানী \L ukasiewicz যুক্তিবিদ্যাও একইভাবে সংজ্ঞায়িত, শুধু
$V = V_m$ নেওয়া হয়।
\end{defn}

\begin{prop}
  \olref[thr][luk]{def:lukasiewicz} অনুযায়ী সংজ্ঞায়িত
  $\LogLuk[3]$ যুক্তিবিদ্যা এবং \olref{def:lukasiewicz} অনুযায়ী
  সংজ্ঞায়িত $\LogLuk[3]$ যুক্তিবিদ্যা একই।
\end{prop}

\begin{proof}
  \olref[thr][luk]{def:lukasiewicz}-এ দেওয়া সংযোজকগুলির সত্যসারণি
  \olref{def:lukasiewicz}-এর সমীকরণ থেকে পাওয়া নিচের সত্যসারণিগুলির
  সঙ্গে তুলনা করলেই তা দেখা যায়:
  \begin{center}
    \begin{tabular}{c|c}
      $\tf{\lnot}$ & \\
      \hline
      $1$ & $0$ \\
      $1/2$ & $1/2$ \\
      $0$ & $1$
    \end{tabular}
    \quad
    \begin{tabular}{c|ccc}
      $\tf{\land}[\LogLuk[3]]$ & $1$ & $1/2$ & $0$ \\
      \hline
      $1$ & $1$ & $1/2$ & $0$ \\
      $1/2$ & $1/2$ & $1/2$ & $0$\\
      $0$ & $0$ & $0$ & $0$
    \end{tabular}
    \\[2ex]
    \begin{tabular}{c|ccc}
      $\tf{\lor}[\LogLuk[3]]$ & $1$ & $1/2$ & $0$ \\
      \hline
      $1$ & $1$ & $1$ & $1$ \\
      $1/2$ & $1$ & $1/2$ & $1/2$ \\
      $0$ & $1$ & $1/2$ & $0$
    \end{tabular}
    \quad
    \begin{tabular}{c|ccc}
      $\tf{\lif}[\LogLuk[3]]$ & $1$ & $1/2$ & $0$ \\
      \hline
      $1$ & $1$ & $1/2$ & $0$ \\
      $1/2$ & $1$ & $1$ & $1/2$ \\
      $0$ & $1$ & $1$ & $1$
    \end{tabular}
  \end{center}
\end{proof}

\begin{prop}\ollabel{prop:luk-infty-m}
  যদি $\Gamma \Entails[\LogLuk[\infty]] !B$ হয়, তবে সব~$m \ge 2$-এর
  জন্য $\Gamma \Entails[\LogLuk[m]] !B$।
\end{prop}

\begin{proof}
  অনুশীলনী।
\end{proof}

\begin{prob}
  \olref[mvl][inf][luk]{prop:luk-infty-m} প্রমাণ করো।
\end{prob}

সসীম $\Gamma$-এর ক্ষেত্রে এর বিপরীত দাবিটিও সত্য।
% BN-SRC-327 TARGET-CORRECTION-BEGIN
\par\smallskip\noindent\textnormal{\small\textbf{উৎস-সংশোধন (BN-SRC-327):} হিমায়িত উৎসে ``বিপরীত দাবিটিও সত্য'' বললেও $\Gamma$-কে সসীম ধরা হয়নি। অসীম পূর্বধারণাসেটের জন্য দাবিটি ব্যর্থ: $p\oplus q := \lnot p\lif q$ লিখে, $n\ge2$-এর জন্য $D_n$-কে $p_n$-এর $(n-1)$-বার $\oplus$-যোগফল ধরি। $\Gamma$-তে প্রতিটি $(D_n\lif\lnot p_n)\land(\lnot p_n\lif D_n)$ রাখলে মনোনীত মান~$1$ পাওয়ার জন্য $p_n=1/n$ চাই। কোনো সসীম $V_m$-ই সব পূর্বধারণা একসঙ্গে পূরণ করতে পারে না (কারণ $1/m\notin V_m$); তাই সেখানে $\Gamma$ থেকে যেকোনো $q$ অনুসিদ্ধ। কিন্তু $V_\infty$-এ সব $p_n=1/n$ এবং $q=0$ রাখলে পূর্বধারণাগুলি সত্য, উপসংহার সত্য নয়। সসীম $\Gamma$ হলে একটি প্রতিবিপরীত মূল্যায়নের সসীমসংখ্যক মূলদ মান কোনো একটি $V_m$-এ একসঙ্গে থাকে; তাই সেই সীমিত ক্ষেত্রে বিপরীত দাবিটি ঠিক।} % BN-SRC-327 TARGET-CORRECTION-END

অসীমমানী \L ukasiewicz যুক্তিবিদ্যা সবচেয়ে জনপ্রিয় ফাজি যুক্তিবিদ্যা।
ফাজি যুক্তিবিদ্যার সাহিত্যে শর্তবচনকে প্রায়ই $\lnot !A \lor !B$
হিসেবে সংজ্ঞায়িত করা হয়। তাতে অসীমমানী শক্তিশালী ক্লিনি যুক্তিবিদ্যা
পাওয়া যাবে।

\begin{prob}
  দেখাও যে $(p \lif q) \lor (q \lif p)$ সূত্রটি
  $\LogLuk[\infty]$-এ সর্বতঃসত্য।
\end{prob}

\end{document}
